% Partie : logique intuitionniste
% Chapitre : sémantique
% Section : propositions

\documentclass[../../../include/open-logic-section]{subfiles}

\begin{document}

\olfileid[fr]{int}{sem}{pro}

\olsection{Propositions}

Il est parfois utile de donner un nom à l'ensemble des mondes
d'!!a{relational model}~$\mModel{M}$ dans lesquels
!!a{formula}~$!A$ est vraie. Nous l'appelons la
\emph{proposition} définie par~$!A$ et le notons
$\Prop{M}{!A}$.

\begin{defn}
Nous prolongeons la fonction $V$ en une application
$\Frm\to\Pow{W},\qquad !A\mapsto\Prop{M}{!A}$
définie inductivement comme suit :\footnote{Précision de
l'édition française : la source emploie une valeur de cette
application à la place de son nom dans la déclaration
de type. Nous distinguons ici l'application de ses valeurs.}
\begin{enumerate}
\item $\Prop{M}{\lfalse}=\emptyset$.
\item $\Prop{M}{p}=V(p)$.
\item $\Prop{M}{\lnot !B}={}$ :
\[
\Setabs{w \in W}{\text{pour tout $w'$ tel que $Rww'$, $w'
    \notin \Prop{M}{!B}$}}.
\]
\item $\Prop{M}{!B\land !C}=\Prop{M}{!B}\cap\Prop{M}{!C}$.
\item $\Prop{M}{!B\lor !C}=\Prop{M}{!B}\cup\Prop{M}{!C}$.
\item $\Prop{M}{!B\lif !C}={}$ :
\[
\Setabs{w \in W}{\text{pour tout $w'$ tel que $Rww'$, si $w'
    \in \Prop{M}{!B}$, alors $w' \in \Prop{M}{!C}$}}.
\]
\end{enumerate}
\end{defn}

\begin{prop}
\ollabel{prop:propositions-true}
$\mSat{M}{!A}[w]$ si et seulement si
$w\in\Prop{M}{!A}$.\footnote{Correction de l'édition
française : la source utilise ici la macro désignant le
modèle au lieu de celle exprimant la vérité d'une formule
dans ce modèle à un monde donné.}
\end{prop}

\begin{proof}
Exercice.
\end{proof}

\begin{prob}
Démontrer la \olref[int][sem][pro]{prop:propositions-true}.
\end{prob}

\end{document}
